% !TEX program = xelatex
\documentclass[11pt,a4paper]{article}

\usepackage{fontspec}
\usepackage{polyglossia}
\usepackage{amsmath,amssymb,amsfonts}
\usepackage{xcolor}
\usepackage{graphicx}
\usepackage{geometry}
\usepackage{booktabs}

\geometry{left=2cm,right=2cm,top=2cm,bottom=2cm}

\setdefaultlanguage{persian}
\setotherlanguage{english}
\setmainfont{Vazir}

\newcommand{\en}[1]{\textenglish{#1}}
\newcommand{\separator}{\vspace{1em}\noindent\rule{\linewidth}{0.4pt}\vspace{1em}}
\newcommand{\thickseparator}{\vspace{1em}\noindent\rule{\linewidth}{1pt}\vspace{1em}}

\title{\textbf{پاسخ سوال ۷} \\[0.4cm] \footnotesize محاسبه \en{Linear Interpolation} بین مدل بای‌گرام و یونی‌گرم}
\author{محمد باغبان زاده \\ شماره دانشجویی: شناسهٔ دانشجو}
\date{}

\linespread{1.35}

\begin{document}

\maketitle
\thickseparator

\section{مفاهیم پایه: \en{Linear Interpolation}}
بر اساس بخش ۳.۶.۳ کتاب، در روش \en{linear interpolation}، تخمین نهایی احتمال یک n-gram به صورت ترکیب وزنی از Probabilities سطوح مختلف (یونی‌گرم، بای‌گرام، ترای‌گرم و ...) محاسبه می‌شود. در این سوال، از ترکیب مدل بای‌گرام و یونی‌گرم با ضرایب مساوی استفاده می‌کنیم:

\[
\hat{P}(w_n|w_{n-1}) = \lambda_1 P(w_n) + \lambda_2 P(w_n|w_{n-1})
\]

که در آن $\lambda_1 = \lambda_2 = \frac{1}{2}$.

\separator
\section{داده‌های مسأله}
پیکره شامل چهار جمله زیر است (جمله سوم تکراری است):

\begin{center}
\begin{tabular}{l}
\en{<s> I am Sam </s>} \\
\en{<s> Sam I am </s>} \\
\en{<s> I am Sam </s>} \\
\en{<s> I do not like green eggs and Sam </s>}
\end{tabular}
\end{center}

\textbf{نکته مهم}: بر اساس صورت سوال، نمادهای \en{<s>} و \en{</s>} دقیقاً مانند سایر توکن‌ها در شمارش‌ها لحاظ می‌شوند.

\separator
\newpage
\section{محاسبه شمارش‌های مورد نیاز}

\subsection*{شمارش بای‌گرام \en{am Sam}}
جملاتی که شامل \en{am Sam} هستند:
\begin{itemize}
\item جمله ۱: \en{<s> I am Sam </s>} → یک بار
\item جمله ۲: \en{<s> Sam I am </s>} → ندارد (\en{am} قبل از \en{Sam} نیست)
\item جمله ۳: \en{<s> I am Sam </s>} → یک بار
\item جمله ۴: \en{<s> I do not like green eggs and Sam </s>} → ندارد (\en{am} وجود ندارد)
\end{itemize}
\[
C(\text{am Sam}) = 2
\]

\subsection*{شمارش یونی‌گرم \en{Sam}}
تعداد دفعات حضور کلمه \en{Sam} در کل پیکره:
\begin{itemize}
\item جمله ۱: یک \en{Sam}
\item جمله ۲: یک \en{Sam}
\item جمله ۳: یک \en{Sam}
\item جمله ۴: یک \en{Sam} (در انتها)
\end{itemize}
\[
C(\text{Sam}) = 4
\]

\subsection*{شمارش یونی‌گرم \en{am}}
تعداد دفعات حضور کلمه \en{am} در کل پیکره:
\begin{itemize}
\item جمله ۱: یک \en{am}
\item جمله ۲: یک \en{am}
\item جمله ۳: یک \en{am}
\item جمله ۴: \en{am} وجود ندارد
\end{itemize}
\[
C(\text{am}) = 3
\]

\subsection*{شمارش کل توکن‌ها (برای یونی‌گرم)}
برای محاسبه $P(\text{Sam})$، باید کل تعداد توکن‌های پیکره (شامل \en{<s>} و \en{</s>}) را بشماریم:
\begin{itemize}
\item جمله ۱: \en{<s>, I, am, Sam, </s>} → ۵ توکن
\item جمله ۲: \en{<s>, Sam, I, am, </s>} → ۵ توکن
\item جمله ۳: \en{<s>, I, am, Sam, </s>} → ۵ توکن
\item جمله ۴: \en{<s>, I, do, not, like, green, eggs, and, Sam, </s>} → ۱۰ توکن
\end{itemize}
\[
N = 5 + 5 + 5 + 10 = 25
\]

\separator
\section{محاسبه احتمالات \en{Maximum Likelihood Estimation (MLE)}}

\subsection*{احتمال بای‌گرام $P(\text{Sam}|\text{am})$}
\[
P_{\text{MLE}}(\text{Sam}|\text{am}) = \frac{C(\text{am Sam})}{C(\text{am})} = \frac{2}{3} \approx 0.6667
\]

\subsection*{احتمال یونی‌گرم $P(\text{Sam})$}
\[
P_{\text{MLE}}(\text{Sam}) = \frac{C(\text{Sam})}{N} = \frac{4}{25} = 0.16
\]

\separator
\section{محاسبه \en{Linear Interpolation}}
با استفاده از فرمول درون‌یابی خطی با $\lambda_1 = \lambda_2 = \frac{1}{2}$:

\[
\begin{aligned}
\hat{P}(\text{Sam}|\text{am}) &= \lambda_1 P(\text{Sam}) + \lambda_2 P(\text{Sam}|\text{am}) \\
&= \frac{1}{2} \times 0.16 + \frac{1}{2} \times 0.6667 \\
&= 0.08 + 0.33335 \\
&= 0.41335
\end{aligned}
\]

\separator
\section{جدول خلاصه محاسبات}
\begin{center}
\begin{tabular}{|l|c|}
\hline
\textbf{مقدار} & \textbf{عدد} \\
\hline
$C(\text{am Sam})$ & ۲ \\
\hline
$C(\text{am})$ & ۳ \\
\hline
$C(\text{Sam})$ & ۴ \\
\hline
$N$ (کل توکن‌ها) & ۲۵ \\
\hline
$P_{\text{MLE}}(\text{Sam}|\text{am})$ & $2/3 \approx 0.6667$ \\
\hline
$P_{\text{MLE}}(\text{Sam})$ & $4/25 = 0.16$ \\
\hline
$\lambda_1 = \lambda_2$ & $1/2 = 0.5$ \\
\hline
$\hat{P}(\text{Sam}|\text{am})$ & $0.08 + 0.33335 = 0.41335$ \\
\hline
\end{tabular}
\end{center}

\separator
\section{نتیجه‌گیری نهایی}
\[
\boxed{\hat{P}(\text{Sam}|\text{am}) = 0.41335 \approx \frac{41335}{100000}}
\]

به عبارت دیگر، احتمال اینکه بعد از کلمه \en{am} کلمه \en{Sam} بیاید، با استفاده از درون‌یابی خطی بین مدل بای‌گرام و یونی‌گرم با وزن‌های مساوی، تقریباً \en{۰.۴۱۳۳۵} محاسبه می‌شود.

\separator
\section{تحلیل نتایج}
\begin{itemize}
\item مدل بای‌گرام به تنهایی احتمال $0.6667$ را تخمین می‌زد (چون از ۳ بار دیده شدن \en{am}، ۲ بار با \en{Sam} همراه بوده است).
\item مدل یونی‌گرم به تنهایی احتمال $0.16$ را تخمین می‌زند (چون \en{Sam} از کل ۲۵ توکن، ۴ بار دیده شده است).
\item ترکیب خطی این دو با وزن مساوی، عددی بین این دو مقدار تولید می‌کند: $0.41335$.
\item این روش به مدل کمک می‌کند تا برای \en{n-gram}هایی که کمتر دیده شده‌اند، تخمین نرم‌تری داشته باشد و احتمال صفر ندهد.
\end{itemize}

\thickseparator

\end{document}